% This file was assembled from main.tex and manuscript inputs.
% Generated on 2026-06-01.

\documentclass{amsart}

%\usepackage[a4paper,margin=30mm]{geometry}
\usepackage{amsmath,amssymb}
\usepackage[colorlinks=true,citecolor=blue,urlcolor=blue]{hyperref}

\theoremstyle{plain}
\newtheorem{theorem}{Theorem}[section]
\newtheorem{lemma}[theorem]{Lemma}
\newtheorem{corollary}[theorem]{Corollary}

\theoremstyle{definition}
\newtheorem{definition}[theorem]{Definition}

\newcommand{\R}{\mathbb{R}}
\newcommand{\Pc}{\mathcal{P}}
\newcommand{\ZFC}{\mathsf{ZFC}}
\newcommand{\FMEA}{\mathsf{FMEA}}
\newcommand{\CH}{\mathsf{CH}}
\newcommand{\Con}{\operatorname{Con}}

\title{Relative Independence of Erd\H{o}s Problem \#501}
\author{Sungchul Lee}
\date{\today}

\begin{document}

\begin{abstract}
We prove a strengthening of the positive assertion in Erd\H{o}s problem \#501 assuming that Lebesgue measure extends to a measure on all subsets of $\R$.
Together with Hechler's counterexample under the continuum hypothesis, this gives a relative independence result for Erd\H{o}s problem \#501, relative to the consistency of this extension hypothesis, and hence relative to the consistency of a measurable cardinal.
\end{abstract}

\maketitle

\section{Introduction}

Let $m$ denote Lebesgue measure on $\R$, and let $m^*$ denote Lebesgue outer measure.
Erd\H{o}s problem \#501 \cite{ErdosProblems501} asks whether every family $(A_x)_{x\in\R}$ of bounded subsets of $\R$ with $m^*(A_x)<1$ admits an infinite independent set, that is, an infinite set $X\subset\R$ such that $x\notin A_y$ whenever $x,y\in X$ are distinct.
We write $P$ for the affirmative answer to this question.

Erd\H{o}s and Hajnal \cite{ErdosHajnal1960} proved the existence of arbitrarily large finite independent sets.
Hechler \cite{Hechler1972} proved $\neg P$ under the assumption $\CH$.
Newelski, Pawlikowski, and Seredy\'nski \cite{NPS1987} proved that if all sets $A_y$ are closed and have measure $<1$, then there is an infinite independent set.
However, it remained open whether $\neg P$ can be proved without assuming $\CH$.

Let $\FMEA$ denote the Full Measure Extension Axiom, the assertion that Lebesgue measure has an extension to a countably additive measure on all subsets of $\R$.
Our main result is the following.

\begin{theorem}\label{thm:main}
Assume $\ZFC+\FMEA$.
Every family $(A_y)_{y\in\R}$ of subsets of $\R$ with Lebesgue outer measure $<1$ admits an infinite independent set.
\end{theorem}

Note that Theorem \ref{thm:main} does not assume boundedness.
Hence Theorem \ref{thm:main} implies $P$ under $\FMEA$.
$\FMEA$ is equiconsistent with the existence of a measurable cardinal; see \cite[1D(e), 2E]{FremlinRVM}.
Combining Theorem \ref{thm:main} with Hechler's theorem gives the following consequence.

\begin{corollary}\label{cor:relative-independence}
If $\Con(\ZFC+\FMEA)$ holds, then $P$ is independent of $\ZFC$.
In particular, if $\ZFC$ with a measurable cardinal is consistent, then $P$ is independent of $\ZFC$.
\end{corollary}

The key ingredient of the proof is Lemma \ref{lem:direct-section-inequality}, which gives a Fubini-type estimate for arbitrary subsets of $\R^2$.
This makes it possible to handle arbitrary sets $A_y$.

\subsection*{Notation}

For a set $X$, we write $\Pc(X)$ to denote its power set.
$\R^{<\omega}$ denotes the set of all finite sequences of real numbers; a sequence $s\in\R^{<\omega}$ of length $n$ is a function from $\{0,\ldots,n-1\}$ to $\R$.

\section{Proof of the Main Theorem}\label{sec:main-theorem-proof}

This section proves Theorem \ref{thm:main} from Lemma \ref{lem:key-selection}.
We first state Lemma \ref{lem:key-selection}, postponing its proof to a later section, and then show how Theorem \ref{thm:main} follows from it.

If $\nu:\Pc(\R)\to[0,\infty]$ is a measure extending Lebesgue measure, then every set $S\subset\R$ satisfies
\begin{equation}\label{eq:outer-measure-domination}
\nu(S)\le m^*(S).
\end{equation}
Indeed, this follows by covering $S$ by countably many open intervals and taking the infimum over all such covers.
In particular, if $m^*(S)<1$, then $\nu(S)<1$.

For a family $(A_y)_{y\in\R}$ of subsets of $\R$, write $B_x=\{y\in\R:x\in A_y\}$ for $x\in\R$.

\begin{lemma}\label{lem:key-selection}
Let $\nu:\Pc(\R)\to[0,\infty]$ be a measure extending Lebesgue measure.
Let $(A_y)_{y\in\R}$ be a family of subsets of $\R$ such that $m^*(A_y)<1$ for every $y\in\R$.
If $C\subset\R$ and $\nu(C)=\infty$, then there is an $x\in C$ such that $\nu(C\setminus B_x)=\infty$.
\end{lemma}

\begin{proof}[Proof of Theorem \ref{thm:main}]
Assume $\FMEA$, and fix a measure $\nu:\Pc(\R)\to[0,\infty]$ extending Lebesgue measure.
Let $(A_y)_{y\in\R}$ satisfy $m^*(A_y)<1$ for every $y\in\R$.
Fix a well-ordering $\preceq$ of $\R$ and a point $r_0\in\R$.

For a finite sequence $s\in\R^{<\omega}$, define
\begin{equation}\label{eq:definition-of-Cs}
C_s
=
\R\setminus\bigcup_{i<|s|}\left(A_{s(i)}\cup B_{s(i)}\cup\{s(i)\}\right),
\end{equation}
where $|s|$ denotes the length of $s$.
Let
\begin{equation}\label{eq:definition-of-Qs}
Q_s
=
\{a\in C_s:\nu(C_s\setminus B_a)=\infty\}.
\end{equation}
Define $G:\R^{<\omega}\to\R$ by
\[
G(s)=
\begin{cases}
\text{the $\preceq$-least element of $Q_s$}, & Q_s\ne\varnothing,\\
r_0, & Q_s=\varnothing.
\end{cases}
\]
We write $f\upharpoonright n$ for the restriction of $f$ to $\{0,\ldots,n-1\}$.
By recursion on $\omega$, there exists $f:\omega\to\R$ such that
\[
f(n)=G(f\upharpoonright n)
\qquad(n<\omega).
\]

We prove by induction that
\[
\nu(C_{f\upharpoonright n})=\infty
\qquad(n<\omega).
\]
For $n=0$, $f\upharpoonright 0$ is the empty sequence, so $C_{f\upharpoonright 0}=\R$ and $\nu(C_{f\upharpoonright 0})=\infty$.
Suppose $\nu(C_{f\upharpoonright n})=\infty$.
Lemma \ref{lem:key-selection} applied to $C_{f\upharpoonright n}$ and \eqref{eq:definition-of-Qs} show that $Q_{f\upharpoonright n}$ is nonempty.
Hence $f(n)=G(f\upharpoonright n)$ belongs to $C_{f\upharpoonright n}$ and satisfies
\[
\nu(C_{f\upharpoonright n}\setminus B_{f(n)})=\infty.
\]
Moreover, by \eqref{eq:definition-of-Cs},
\[
C_{f\upharpoonright(n+1)}
=
C_{f\upharpoonright n}\setminus\left(A_{f(n)}\cup B_{f(n)}\cup\{f(n)\}\right)
=
(C_{f\upharpoonright n}\setminus B_{f(n)})\setminus\left(A_{f(n)}\cup\{f(n)\}\right).
\]
By \eqref{eq:outer-measure-domination},
\[
\nu(A_{f(n)})\le m^*(A_{f(n)})<1,
\]
and $\nu(\{f(n)\})=0$ because $\nu$ extends Lebesgue measure.
Thus $\nu(A_{f(n)}\cup\{f(n)\})<\infty$.
Since $\nu(C_{f\upharpoonright n}\setminus B_{f(n)})=\infty$ and $\nu(A_{f(n)}\cup\{f(n)\})<\infty$, we have $\nu(C_{f\upharpoonright(n+1)})=\infty$.
This completes the induction.

Set $X=\{f(n):n<\omega\}$.
Let $i<j$.
By \eqref{eq:definition-of-Cs},
\[
f(j)\in C_{f\upharpoonright j}\subset C_{f\upharpoonright(i+1)}
=C_{f\upharpoonright i}\setminus\left(A_{f(i)}\cup B_{f(i)}\cup\{f(i)\}\right),
\]
we have $f(j)\ne f(i)$, $f(j)\notin A_{f(i)}$, and $f(j)\notin B_{f(i)}$.
The last relation is equivalent to $f(i)\notin A_{f(j)}$.
Thus $X$ is infinite, and $x\notin A_y$ for all distinct $x,y\in X$.
\end{proof}

\section{Proof of Lemma \ref{lem:key-selection}}\label{sec:key-selection}

We use the following elementary section inequality.
Here $\overline{\int_\R} g\,dm$ denotes the Lebesgue upper integral; that is, for $g:\R\to[0,\infty]$,
\begin{equation}\label{eq:upper-integral-definition}
\overline{\int_\R} g\,dm
=
\inf\left\{\int_\R h\,dm:g\le h,\ h\text{ is Lebesgue measurable}\right\}.
\end{equation}

\begin{lemma}\label{lem:direct-section-inequality}
Let $(Y,\Pc(Y),\nu)$ be a $\sigma$-finite measure space.
For every set $H\subset\R\times Y$,
\begin{equation}\label{eq:direct-section-inequality}
\overline{\int_\R}\nu(H_x)\,dm(x)
\le
\int_Y m^*(H^y)\,d\nu(y),
\end{equation}
where
\begin{equation}\label{eq:section-notation}
H_x=\{y\in Y:(x,y)\in H\},
\qquad
H^y=\{x\in\R:(x,y)\in H\}.
\end{equation}
\end{lemma}

\begin{proof}
Since $\nu$ is $\sigma$-finite, there is a $\nu$-measurable function $\eta:Y\to(0,\infty)$ such that
\begin{equation}\label{eq:eta-integral-bound}
\int_Y \eta\,d\nu
\le
1.
\end{equation}
Fix $\varepsilon>0$.
For each $y\in Y$, choose an open set $U_y\subset\R$ such that
\begin{equation}\label{eq:choice-of-Uy}
H^y\subset U_y
\qquad\text{and}\qquad
m(U_y)\le m^*(H^y)+\varepsilon\eta(y).
\end{equation}
If $m^*(H^y)=\infty$, we may simply take $U_y=\R$.

Let $(I_n)_{n<\omega}$ be a countable base for the usual topology of $\R$, for instance the open intervals with rational endpoints.
For each $n<\omega$, define
\[
Y_n=\{y\in Y:I_n\subset U_y\}.
\]
Each $Y_n$ is $\nu$-measurable, because the $\sigma$-algebra on $Y$ is $\Pc(Y)$.
Set
\[
E=\bigcup_{n<\omega}(I_n\times Y_n).
\]
Then $E$ is measurable in the product of the Lebesgue $\sigma$-algebra on $\R$ and $\Pc(Y)$.

We claim that $E^y=U_y$ for every $y\in Y$.
Indeed, if $x\in E^y$, then $x\in I_n$ for some $n$ with $I_n\subset U_y$, so $x\in U_y$.
Conversely, if $x\in U_y$, then, since $U_y$ is open, there is a basic interval $I_n$ such that $x\in I_n\subset U_y$.
Thus $y\in Y_n$, and hence $x\in E^y$.

Therefore $H\subset E$.
By Tonelli's theorem applied to the measurable set $E$,
\begin{equation}\label{eq:tonelli-for-E}
\int_\R \nu(E_x)\,dm(x)
=
\int_Y m(E^y)\,d\nu(y)
=
\int_Y m(U_y)\,d\nu(y).
\end{equation}
By \eqref{eq:eta-integral-bound} and \eqref{eq:choice-of-Uy},
\begin{equation}\label{eq:Uy-integral-bound}
\int_Y m(U_y)\,d\nu(y)
\le
\int_Y m^*(H^y)\,d\nu(y)+\varepsilon\int_Y \eta\,d\nu
\le
\int_Y m^*(H^y)\,d\nu(y)+\varepsilon.
\end{equation}
Since $H_x\subset E_x$ for every $x\in\R$, we have $\nu(H_x)\le\nu(E_x)$.
Moreover, $x\mapsto\nu(E_x)$ is Lebesgue measurable.
Hence, by \eqref{eq:upper-integral-definition}, \eqref{eq:tonelli-for-E}, and \eqref{eq:Uy-integral-bound},
\[
\overline{\int_\R}\nu(H_x)\,dm(x)
\le
\int_\R \nu(E_x)\,dm(x)
\le
\int_Y m^*(H^y)\,d\nu(y)+\varepsilon.
\]
Letting $\varepsilon\to0$ gives \eqref{eq:direct-section-inequality}.
\end{proof}

For the measure $\nu$ in Lemma \ref{lem:key-selection}, $(\R,\Pc(\R),\nu)$ is $\sigma$-finite, since
\[
\R=\bigcup_{n\ge1}[-n,n]
\qquad\text{and}\qquad
\nu([-n,n])=2n<\infty.
\]
Applying Lemma \ref{lem:direct-section-inequality} with $Y=\R$, we get, for every $H\subset\R^2$,
\begin{equation}\label{eq:section-inequality}
\overline{\int_\R}\nu(H_x)\,dm(x)
\le
\int_\R m^*(H^y)\,d\nu(y),
\end{equation}
where $H_x$ and $H^y$ are as in \eqref{eq:section-notation}, with $Y=\R$.

\begin{proof}[Proof of Lemma \ref{lem:key-selection}]
Let $C\subset\R$ satisfy $\nu(C)=\infty$.
Suppose, toward a contradiction, that
\begin{equation}\label{eq:contrary}
\nu(C\setminus B_x)<\infty,
\qquad
\forall x\in C.
\end{equation}

For $N\ge1$, put $C_N=C\cap[-N,N]$.
Since $C=\bigcup_{N\ge1}C_N$ and the sequence $(C_N)_{N\ge1}$ is increasing, continuity from below gives
\begin{equation}\label{eq:CN-to-infinity}
\nu(C_N)\longrightarrow\infty.
\end{equation}
Choose $N$ such that $D:=C\cap[-N,N]$ satisfies $\nu(D)>1$.
For $k\ge1$, define
\begin{equation}\label{eq:definition-of-Dk}
D_k=\{x\in D:\nu(C\setminus B_x)\le k\}.
\end{equation}
By \eqref{eq:contrary}, $D=\bigcup_{k\ge1}D_k$, and the sequence $(D_k)_{k\ge1}$ is increasing.
Again by continuity from below, there is $k\ge1$ such that $\nu(D_k)>1$.
By \eqref{eq:outer-measure-domination} and $D_k\subset[-N,N]$,
\begin{equation}\label{eq:Dk-outer-bounds}
1<m^*(D_k)<\infty.
\end{equation}

For every integer $M>N$,
\[
\nu(C_M)\le\nu([-M,M])=2M<\infty.
\]
By \eqref{eq:CN-to-infinity} and \eqref{eq:Dk-outer-bounds}, choose an integer $M>N$ such that
\begin{equation}\label{eq:choice-of-M}
\nu(C_M)>\frac{k\,m^*(D_k)}{m^*(D_k)-1}.
\end{equation}
In particular, $\nu(C_M)>k$.

If $x\in D_k$, then $\nu(C\setminus B_x)\le k$ by \eqref{eq:definition-of-Dk}.
Since $C_M$ has finite $\nu$-measure and $C_M\setminus B_x\subset C\setminus B_x$, we have
\begin{equation}\label{eq:vertical-lower-section}
\nu(B_x\cap C_M)
=
\nu(C_M)-\nu(C_M\setminus B_x)
\ge
\nu(C_M)-k,
\qquad
\forall x\in D_k.
\end{equation}

Let
\[
H=\{(x,y)\in D_k\times C_M:x\in A_y\}\subset\R^2.
\]
For $x\in\R$,
\[
H_x=
\begin{cases}
B_x\cap C_M, & x\in D_k,\\
\varnothing, & x\notin D_k.
\end{cases}
\]
Thus \eqref{eq:vertical-lower-section} implies
\[
\nu(H_x)\ge(\nu(C_M)-k)1_{D_k}(x)
\qquad(x\in\R).
\]
By monotonicity of upper integral and the identity
\[
\overline{\int_\R} a1_S\,dm=a\,m^*(S)
\qquad(a\ge0,
\ S\subset\R),
\]
we get
\begin{equation}\label{eq:left-lower}
\overline{\int_\R}\nu(H_x)\,dm(x)
\ge
(\nu(C_M)-k)m^*(D_k).
\end{equation}

For $y\in\R$,
\[
H^y=
\begin{cases}
A_y\cap D_k, & y\in C_M,\\
\varnothing, & y\notin C_M.
\end{cases}
\]
Since $m^*(A_y)<1$ for every $y$,
\[
m^*(H^y)
\le
1_{C_M}(y)
\qquad(y\in\R).
\]
The function $y\mapsto m^*(H^y)$ is $\nu$-measurable because $\nu$ is defined on all subsets of $\R$.
Hence
\begin{equation}\label{eq:right-upper}
\int_\R m^*(H^y)\,d\nu(y)
\le
\int_\R 1_{C_M}(y)\,d\nu(y)
=
\nu(C_M).
\end{equation}

Applying \eqref{eq:section-inequality} to this $H$ and using \eqref{eq:left-lower} and \eqref{eq:right-upper}, we obtain
\[
(\nu(C_M)-k)m^*(D_k)
\le
\nu(C_M).
\]
Dividing by $\nu(C_M)>0$ gives
\[
\left(1-\frac{k}{\nu(C_M)}\right)m^*(D_k)
\le
1.
\]
However, \eqref{eq:choice-of-M} gives
\[
\left(1-\frac{k}{\nu(C_M)}\right)m^*(D_k)>1,
\]
contradicting the preceding inequality.
Therefore \eqref{eq:contrary} is false, and some $x\in C$ satisfies $\nu(C\setminus B_x)=\infty$.
\end{proof}

\appendix

\section{Counterexample Under CH}\label{sec:ch-counterexample}

For completeness, we give a simple proof that $P$ fails under $\CH$.
Assume $\CH$, and fix an enumeration $\R=\{r_\alpha:\alpha<\omega_1\}$.
Let $\preceq$ be the induced well-ordering of $\R$, so that $r_\alpha\preceq r_\beta$ if and only if $\alpha\le\beta$, and let $\prec$ be the corresponding strict order.
For every $\beta<\omega_1$, the set $\{r_\alpha:\alpha<\beta\}$ is countable.
For $y=r_\beta$, define
\[
A_y
=
\{r_\alpha:\alpha<\beta,\ |r_\alpha|\le |y|+1\}.
\]
Then $A_y$ is countable, and hence $m^*(A_y)=0$.
Also $A_y\subset[-(|y|+1),|y|+1]$, so $A_y$ is bounded.

We claim that the family $(A_y)_{y\in\R}$ has no infinite independent set.
Suppose otherwise, and let $X\subset\R$ be infinite and independent.
Choose a strictly increasing sequence
\[
x_0\prec x_1\prec x_2\prec\cdots
\]
from $X$.
If $i<j$, then independence gives $x_i\notin A_{x_j}$.
Since $x_i\prec x_j$, the definition of $A_{x_j}$ implies
\[
|x_i|>|x_j|+1.
\]
In particular, $|x_n|>|x_{n+1}|+1$ for every $n<\omega$.
Iterating, we get $|x_n|<|x_0|-n$, which is impossible for $n>|x_0|$.
Therefore no infinite independent set exists, and $P$ fails under $\CH$.

\section*{Use of AI}

During the preparation of this work, the author used OpenAI's GPT-5.5 Pro to explore possible proof strategies and to obtain the main methodological ideas behind the argument.
After using this tool, the author independently verified the mathematical details, refined the exposition, and takes full responsibility for the final manuscript.


\begin{thebibliography}{1}

\bibitem{ErdosProblems501}
T.~F. Bloom.
\newblock {Erd\H{o}s Problem \#501}.
\newblock Accessed 29 May 2026.
\newblock URL: \url{https://www.erdosproblems.com/501}.

\bibitem{ErdosHajnal1960}
P.~Erd\H{o}s and A.~Hajnal.
\newblock Some remarks on set theory. {VIII}.
\newblock {\em Michigan Math. J.}, 7:187--191, 1960.

\bibitem{FremlinRVM}
D.~H. Fremlin.
\newblock Real-valued-measurable cardinals.
\newblock Version of 19 September 2009.
\newblock URL: \url{https://www1.essex.ac.uk/maths/people/fremlin/rvmc.pdf}.

\bibitem{Hechler1972}
S.~H. Hechler.
\newblock On two problems in combinatorial set theory.
\newblock {\em Bull. Acad. Polon. Sci. S\'er. Sci. Math. Astronom. Phys.}, 20:429--431, 1972.

\bibitem{NPS1987}
L.~Newelski, J.~Pawlikowski, and W.~Seredy\'nski.
\newblock Infinite free set for small measure set mappings.
\newblock {\em Proc. Amer. Math. Soc.}, 100:335--339, 1987.

\end{thebibliography}


\end{document}
